% Casos de uso de la conducción y el análisis (UML): durante la sesión o después de ella.
% CU10 y CU14 extienden a CU9; el simulador participa en lanzar el juego y en la sesión.
\begin{figure}[H]
	\centering
	\resizebox{\textwidth}{!}{%
		\begin{tikzpicture}[
			>=Stealth,
			cu/.style={draw, thick, ellipse, align=center, minimum width=4.8cm, minimum height=1.15cm, font=\small, fill=white},
			dep/.style={->, dashed, thick},
			estereo/.style={font=\scriptsize, fill=gray!6, inner sep=1pt},
			actor/.pic={
				\draw[thick] (0,0.62) circle (0.13);
				\draw[thick] (0,0.49) -- (0,0.15);
				\draw[thick] (-0.25,0.38) -- (0.25,0.38);
				\draw[thick] (0,0.15) -- (-0.18,-0.15);
				\draw[thick] (0,0.15) -- (0.18,-0.15);
			}
			]
			\draw[thick, fill=gray!6] (2.9,-0.1) rectangle (15.4,9.7);
			\node[font=\small\bfseries, anchor=north] at (9.15,9.6) {VH SimGrid};

			\node[cu] (c8)  at (6.1,8.4) {CU8. Lanzar un juego};
			\node[cu] (c9)  at (6.1,6.8) {CU9. Conducir una\\sesión asistida};
			\node[cu] (c14) at (6.1,4.9) {CU14. Configurar el equipo\\desde la pantalla del aro};
			\node[cu] (c11) at (6.1,3.3) {CU11. Planificar la\\estrategia de carrera};
			\node[cu] (c12) at (6.1,2.0) {CU12. Diseñar un dash};
			\node[cu] (c13) at (6.1,0.7) {CU13. Analizar una\\sesión grabada};
			\node[cu] (c10) at (12.2,5.0) {CU10. Consultar al\\ingeniero por voz};

			\draw[dep] (c10) -- node[estereo, pos=0.45, above, sloped] {\guillemotleft extend\guillemotright} (c9);
			\draw[dep] (c14) -- node[estereo, right] {\guillemotleft extend\guillemotright} (c9);

			\pic at (0.8,4.4) {actor};
			\node[font=\small, below] at (0.8,4.0) {Piloto};
			\pic at (17.6,7.6) {actor};
			\node[font=\small, below, align=center] at (17.6,7.2) {Simulador\\de conducción};

			\foreach \c in {c8,c9,c14,c11,c12,c13} \draw (1.2,4.75) -- (\c.west);
			\draw (17.2,7.95) -- (c8.east);
			\draw (17.2,7.95) -- (c9.east);
	\end{tikzpicture}}
	\caption[Diagrama de casos de uso de la conducción y el análisis]{Diagrama de casos de uso de la conducción y el análisis. (Fuente propia)}
	\label{fig:cu-conduccion}
\end{figure}
